\documentclass[11pt]{article}

\usepackage[a4paper,margin=1in]{geometry}
\usepackage{amsmath,amssymb,amsthm}
\usepackage{mathtools}
\usepackage{array}
\usepackage{booktabs}
\usepackage{url}
\usepackage{hyperref}

\title{Clifford Algebras and Their Nonassociative Twins:\\
Passive Volume, Active Volume, and Folded Triality}
\author{Mirco A. Mannucci}
\date{\today}

\newtheorem{definition}{Definition}
\newtheorem{proposition}{Proposition}
\newtheorem{remark}{Remark}
\newtheorem{problem}{Problem}

\newcommand{\Ztwo}{\mathbb Z_2}
\newcommand{\R}{\mathbb R}
\newcommand{\Cl}{\mathrm{Cl}}
\newcommand{\OO}{\mathbb O}
\newcommand{\assoc}{\Phi}
\newcommand{\eps}{\varepsilon}

\begin{document}
\maketitle

\begin{abstract}
Clifford algebras are associative, while the octonions are nonassociative.
The two structures nevertheless meet on the same eight-dimensional linear
body in the case of $\Cl(3,0)$. The key observation is not that octonions are
hidden inside the ordinary Clifford product. Rather, the intrinsic split
\[
\Cl(3,0)=\Cl^+(3,0)\oplus \Cl^+(3,0)I,
\qquad
\Cl^+(3,0)\cong\mathbb H,
\]
is already the Cayley--Dickson doubling space. The standard Clifford product
uses the pseudoscalar $I$ as a passive central volume imaginary. The octonion
product is obtained by keeping the quaternionic rotor core and the same
vector-space split, but replacing the central rule $Ia=aI$ by the conjugating
rule $\ell a=\widetilde a\ell$. Thus the octonionic twin of $\Cl(3,0)$ is
read as an activation of the oriented volume line: passive background
orientation becomes an active chirality-reversing gate.

This note develops that interpretation in three layers. First, it records the
standard twisted Boolean-cube description, where Clifford algebras and
octonions differ by whether the associator is trivial or curved. Second, it
replaces the basis-level sign picture by an invariant Clifford-body picture:
one should not keep only the naked algebra, nor only a presentation
$\Cl(V,Q)$, but the graded body $(\mathcal C,\alpha,\widetilde{(\cdot)},\mathfrak v)$
with reversion and central volume line. Third, it connects the construction to
triality. The usual $\Cl(8,0)$ approach presents triality in unfolded form, as
three distinct eight-dimensional representations in a large associative
envelope. The octonionic approach presents triality in folded form, compressing
the vector and two spinor roles into one nonassociative eight-dimensional body.
The contribution claimed here is not a new construction of the octonions, but
a geometric reframing of the Cayley--Dickson step inside $\Cl(3,0)$ and of the
octonions as a compact folded triality cell.
\end{abstract}

\section{The guiding intuition}

Noncommutativity says that order matters:
\[
xy \neq yx.
\]
Nonassociativity says that parenthesization, or history of composition,
matters:
\[
(xy)z \neq x(yz).
\]
This distinction is already visible in elementary geometry. If
\[
x\star y = \frac{x+y}{2}
\]
is the midpoint operation, then
\[
(x\star y)\star z = \frac{x+y+2z}{4},
\qquad
x\star(y\star z) = \frac{2x+y+z}{4}.
\]
Thus staged aggregation remembers its history.

The Clifford--octonion relation considered here gives a sharper and more
geometric version of the same idea. The octonionic product on the
$\Cl(3,0)$ vector space is not obtained by merely transporting an arbitrary
multiplication table. It is obtained by changing the role of the intrinsic
volume slot. Under the Clifford product, the volume element is passive and
central. Under the octonionic product, the same slot becomes active and
conjugating. The resulting nonassociativity records where this active volume
operation occurs inside a sequence of rotations.

\section{Twisted hypercube algebras}

Let
\[
G = (\Ztwo)^n
\]
be the Boolean cube, written additively. Let $A_F$ be the real vector space
with basis
\[
\{u_a : a\in G\}.
\]
Given a normalized sign twist
\[
F:G\times G\to \{\pm 1\},
\qquad
F(0,a)=F(a,0)=1,
\]
define a product by
\[
u_a\cdot_F u_b = F(a,b)u_{a+b}.
\]
This is a signed deformation of the group algebra $\R[G]$.

\begin{definition}
The associator factor of $F$ is the function
\[
\assoc_F(a,b,c)
=
\frac{F(a,b)F(a+b,c)}
     {F(b,c)F(a,b+c)}.
\]
\end{definition}

A direct computation gives
\[
(u_a\cdot_F u_b)\cdot_F u_c
=
F(a,b)F(a+b,c)u_{a+b+c},
\]
while
\[
u_a\cdot_F (u_b\cdot_F u_c)
=
F(b,c)F(a,b+c)u_{a+b+c}.
\]
Therefore
\[
(u_a\cdot_F u_b)\cdot_F u_c
=
\assoc_F(a,b,c)\,
u_a\cdot_F (u_b\cdot_F u_c).
\]

\begin{proposition}
The twisted hypercube algebra $A_F$ is associative if and only if
\[
\assoc_F(a,b,c)=1
\]
for all $a,b,c\in G$.
Equivalently, $F$ is a $2$-cocycle.
\end{proposition}

This simple formula is the whole mechanism. Clifford algebras are obtained
when the twist has trivial associator. Octonionic-type algebras arise when
the same hypercube skeleton is equipped with a nontrivial but controlled
associator.

\section{Clifford algebras as associative twisted hypercubes}

The real Clifford algebra $\Cl(p,q)$ is generated by elements
\[
e_1,\dots,e_n,\qquad n=p+q,
\]
with relations
\[
e_i^2 =
\begin{cases}
+1, & 1\leq i\leq p,\\
-1, & p<i\leq n,
\end{cases}
\qquad
e_i e_j = - e_j e_i
\quad (i\neq j).
\]
Its standard basis is indexed by subsets of $\{1,\dots,n\}$, equivalently by
elements of $(\Ztwo)^n$:
\[
e_a = e_1^{a_1}\cdots e_n^{a_n},
\qquad
a=(a_1,\dots,a_n)\in(\Ztwo)^n.
\]
Thus
\[
\dim_\R \Cl(p,q)=2^n.
\]

Multiplication takes the form
\[
e_a e_b = F_{\Cl}(a,b)e_{a+b}
\]
for a specific sign twist $F_{\Cl}$. The defining fact is:
\[
\assoc_{F_{\Cl}}(a,b,c)=1.
\]
So Clifford algebras are associative twisted hypercubes.

The commutation signs are still nontrivial. In general,
\[
e_a e_b
=
\beta(a,b) e_b e_a,
\]
where
\[
\beta(a,b)
=
\frac{F_{\Cl}(a,b)}{F_{\Cl}(b,a)}.
\]
Thus Clifford multiplication already contains a controlled sign geometry.
But the associator remains flat.

\section{The octonions as the first nonassociative twin}

The octonions $\OO$ also have dimension $8=2^3$. Their basis may be written
\[
1,u_1,\dots,u_7,
\]
or more structurally as
\[
\{u_a:a\in(\Ztwo)^3\}.
\]
The multiplication again has the form
\[
u_a u_b = F_{\OO}(a,b)u_{a+b},
\]
but now the associator factor
\[
\assoc_{F_{\OO}}(a,b,c)
\]
is nontrivial.

Thus the octonions are not merely ``like'' Clifford algebras because they
have the same dimension as $\Cl(3,0)$. They are like Clifford algebras
because both live on the same $(\Ztwo)^3$-graded hypercube skeleton:
\[
\Cl(3,0): \quad 8\text{-dimensional associative twisted hypercube},
\]
\[
\OO: \quad 8\text{-dimensional nonassociative twisted hypercube}.
\]

The Fano plane is the visible combinatorial shadow of this multiplication.
Its oriented lines encode the products among the seven imaginary basis
elements. The hidden algebraic datum is the nontrivial associator.


\section{The invariant $\Cl(3,0)$ construction: passive volume versus conjugating volume}

The sign-table description is useful, but it is not fully satisfying. It
makes the octonion product look as if it were imposed externally on a chosen
basis. In the special case of $\Cl(3,0)$ there is a more invariant and more
geometric construction.

Let
\[
\Cl(3,0)=\langle e_1,e_2,e_3\mid e_i^2=1,\ e_i e_j=-e_j e_i\rangle.
\]
The even subalgebra is
\[
\Cl^+(3,0)=\langle 1,e_{23},e_{31},e_{12}\rangle \cong \mathbb H.
\]
Geometrically, this is the algebra of three-dimensional rotors. The bivectors
$e_{23},e_{31},e_{12}$ are oriented rotation planes, and unit elements of
$\Cl^+(3,0)$ act as spinors/rotors.

Let
\[
I=e_1e_2e_3.
\]
Then
\[
I^2=-1,
\]
and, in $\Cl(3,0)$, the pseudoscalar $I$ is central. Moreover,
\[
\Cl(3,0)=\Cl^+(3,0)\oplus \Cl^+(3,0)I.
\]
Thus every element can be written uniquely as
\[
a+bI,
\qquad a,b\in \Cl^+(3,0)\cong\mathbb H.
\]
Under the ordinary Clifford product, since $I$ is central, one obtains
\[
(a+bI)(c+dI)=(ac-bd)+(ad+bc)I.
\]
So $\Cl(3,0)$ is a quaternionic rotor algebra doubled by a passive central
volume imaginary.

Now comes the octonionic move. Keep the same vector-space splitting
\[
\Cl^+(3,0)\oplus \Cl^+(3,0)I,
\]
but do not let the second copy be controlled by a central pseudoscalar.
Instead, introduce a new symbol $\ell$ occupying the same second-copy slot and
impose the crossing rule
\[
\ell a=\bar a\,\ell,
\]
where $\bar a$ is quaternionic conjugation on $\Cl^+(3,0)$. In Clifford terms,
this conjugation is simply reversion restricted to the even subalgebra:
\[
\overline{a_0+a_1e_{23}+a_2e_{31}+a_3e_{12}}
=
 a_0-a_1e_{23}-a_2e_{31}-a_3e_{12}.
\]
Then define
\[
(a+b\ell)\star(c+d\ell)
:=
(ac-\bar d\,b)+(da+b\bar c)\ell.
\]
This is the Cayley--Dickson doubling product
\[
\mathbb O\cong \mathbb H\oplus \mathbb H\ell.
\]

Therefore the octonion product is definable from the internal Clifford
structure of $\Cl(3,0)$ together with one additional algebraic operation:
reversion/conjugation on the even quaternionic subalgebra. The construction is
not an arbitrary Fano-plane table placed on the vector space. It is a change in
the role of the pseudoscalar direction.

In ordinary $\Cl(3,0)$, the volume element is passive:
\[
Ia=aI.
\]
In the octonionic twin, the volume-like direction is active:
\[
\ell a=\bar a\ell.
\]
Thus the new direction reverses the orientation of the rotor it crosses. For a
unit rotor $R\in\Cl^+(3,0)$,
\[
\ell R=\bar R\ell=R^{-1}\ell.
\]
So the extra direction behaves like a handedness-reversing or mirror-like gate
for the three-dimensional rotation algebra.

This also explains the nonassociativity. Quaternionic conjugation is an
anti-automorphism:
\[
\overline{ab}=\bar b\bar a.
\]
Therefore crossing the gate before or after composing rotations gives different
answers. Schematically,
\[
\ell(ab)=\overline{ab}\,\ell=\bar b\bar a\ell,
\]
whereas sequentially crossing $a$ and then $b$ gives the opposite order,
\[
(\ell a)b\sim \bar a\bar b\ell.
\]
Since quaternions are noncommutative, generally
\[
\bar a\bar b\neq \bar b\bar a.
\]
Hence parenthesization records where the handedness-reversing gate was inserted into
a sequence of rotations.

The geometric meaning is therefore:
\begin{quote}
$\Cl(3,0)$ is three-dimensional rotor geometry plus a passive volume imaginary.
The octonions are three-dimensional rotor geometry plus an active
handedness-reversing gate.
\end{quote}

Equivalently:
\begin{quote}
$\Cl(3,0)$ doubles $\mathbb H$ by a central $I$; the octonions double
$\mathbb H$ by a conjugating $\ell$.
\end{quote}

This is the invariant version of the ``nonassociative twin'' idea in dimension
$8$. The twisted-hypercube/Fano-plane description records the multiplication
table. The Cayley--Dickson description explains what the operation is doing:
it turns the central volume direction into an orientation-reversing interface.


\section{Is the extra unit determined?}

The preceding construction uses the pseudoscalar direction as the new
Cayley--Dickson unit. This raises an important question: is that choice forced,
or did we merely pick a convenient element of the $8$-dimensional space?

The answer depends on the invariance requirement.

Let
\[
H:=\Cl^+(3,0)\cong \mathbb H,
\qquad
V:=\Cl(3,0)=H\oplus HI.
\]
The odd part of $\Cl(3,0)$ decomposes as
\[
\Cl^-(3,0)=\mathbb R^3\oplus \mathbb R I.
\]
Under the action of $\mathrm{Spin}(3)$, the vector part $\mathbb R^3$ rotates,
whereas the pseudoscalar line $\mathbb R I$ is invariant. Thus, if we demand
that the construction introduce no preferred spatial axis, the only possible
one-dimensional odd direction is
\[
\mathbb R I.
\]
Consequently the additional Cayley--Dickson unit is forced to be
\[
\ell=I
\]
up to sign. The ambiguity
\[
I\mapsto -I
\]
just reverses the chosen orientation and gives an isomorphic octonion algebra.

If we drop rotational invariance, the choice is no longer unique. One can use
an element of the form
\[
qI,
\qquad q\in H,
\qquad q\bar q=1,
\]
as a second-copy generator. But a non-real unit quaternion $q$ encodes a
specific rotor direction or frame. Such a choice is not wrong, but it smuggles
in additional gauge/frame data. It is not the canonical isotropic construction.

So the invariant statement is:
\[
\boxed{
\text{In oriented }\Cl(3,0),\text{ the natural active volume unit is }I=e_1e_2e_3,
\text{ unique up to sign.}
}
\]
This matters conceptually. The octonionic extension is not adding a fourth
spatial direction. It is activating the oriented $3$-volume itself. It is a
chiral extension, not a directional one.

\section{One vector space, two algebra structures}

It is useful to separate the linear object from the multiplication. The real
vector space
\[
V:=\Cl(3,0)
\]
has dimension $8$. But an algebra structure on $V$ is an additional map
\[
\mu:V\otimes V\to V.
\]
The Clifford algebra and the octonions correspond to two different
multiplication maps on the same underlying vector space:
\[
\mu_{\Cl}
\qquad\text{and}\qquad
\mu_{\OO}.
\]
They are not two descriptions of one product. They are two different products
on the same linear body.

The relation is not arbitrary, however. Both products use the same intrinsic
splitting
\[
V=H\oplus HI,
\qquad
H=\Cl^+(3,0)\cong\mathbb H.
\]
For
\[
x=a+bI,
\qquad
c+dI\in H\oplus HI,
\]
the Clifford product is
\[
(a+bI)(c+dI)
=
(ac-bd)+(ad+bc)I.
\]
This is the product in which $I$ is a passive central volume imaginary:
\[
Ia=aI.
\]

The octonionic product keeps the same linear splitting but changes the crossing
law:
\[
Ia=\bar a I.
\]
Equivalently,
\[
(a+bI)\star(c+dI)
=
(ac-\bar d b)+(da+b\bar c)I.
\]
Thus the exact relation between the two structures is:
\[
\boxed{
\text{same quaternionic core, same volume slot, different law for crossing the volume slot.}
}
\]
The Clifford law is central. The octonionic law is conjugating.

This is the clean way to understand how the same $8$-dimensional space can
carry both structures. The vector space is the stage. The product determines
what kind of geometry is being performed on that stage.


\section{From presented Clifford algebras to intrinsic Clifford bodies}

The previous section suggests a more conceptual formulation. In many contexts
one writes a Clifford algebra as
\[
\Cl(V,Q),
\]
where the vector space $V$ and quadratic form $Q$ are treated as primary.
For the present purpose, however, this presentation is not the whole story.
The octonionic twin construction becomes cleaner if we separate three levels:
\[
\text{presented Clifford algebra } \Cl(V,Q),
\]
\[
\text{intrinsic Clifford body } (\mathcal C,\alpha,\widetilde{(\cdot)},\mathfrak v),
\]
\[
\text{choice of product law } \mu.
\]
Here $\alpha$ denotes the grade involution, $\widetilde{(\cdot)}$ denotes
reversion, and $\mathfrak v$ denotes the volume line. The point is not to
forget all structure. The point is to forget the particular generating
presentation only after keeping the internal data needed to recover the even
core, the conjugation operation, and the oriented volume slot.

This distinction matters because the raw associative algebra alone is not
always enough. If we keep only the naked real algebra underlying $\Cl(3,0)$,
then
\[
\Cl(3,0)\cong M_2(\mathbb C)
\]
as a real associative algebra. From this algebra one can recover its center,
\[
Z(\mathcal C)\cong \mathbb C.
\]
The central square roots of $-1$ are precisely the two choices
\[
I \quad\text{and}\quad -I,
\]
so the central complex or volume line
\[
\mathfrak v=\mathbb R I
\]
is intrinsic to the algebra. But the sign of $I$ is not intrinsic: choosing
$I$ rather than $-I$ is exactly the choice of orientation.

Thus the right statement is not that the volume element is absolutely singled
out from nothing. Rather:
\[
\boxed{
\text{The volume line }\mathbb RI\text{ is intrinsic; the volume element }I
\text{ requires orientation.}
}
\]
This mirrors ordinary geometry. An unoriented three-dimensional Euclidean
space has a volume density, but not a signed volume form. An oriented one has
a chosen volume form, up to the sign convention fixed by orientation.

Moreover, the even subalgebra is not recovered merely from the center. It is
recovered from the Clifford body structure, namely from the grade involution:
\[
\mathcal C^+=\operatorname{Fix}(\alpha).
\]
In the case under consideration,
\[
\mathcal C^+=\Cl^+(3,0)\cong\mathbb H.
\]
The reversion $\widetilde{(\cdot)}$ then restricts to quaternionic conjugation
on this even core:
\[
a\mapsto \widetilde a=\bar a,
\qquad a\in \mathcal C^+.
\]
The oriented volume element is characterized internally, up to sign, by the
conditions
\[
I\in Z(\mathcal C),
\qquad
I^2=-1,
\qquad
\alpha(I)=-I,
\qquad
\widetilde I=-I,
\]
and by the decomposition
\[
\mathcal C=\mathcal C^+\oplus \mathcal C^+I.
\]
For $\Cl(3,0)$ this is exactly the decomposition
\[
\mathcal C=H\oplus HI,
\qquad H=\mathcal C^+\cong\mathbb H.
\]

The twin construction is therefore intrinsic in the following precise sense.
Start from an oriented graded Clifford body
\[
(\mathcal C,\alpha,\widetilde{(\cdot)},\mathfrak v),
\]
choose the oriented element $I\in\mathfrak v$ with $I^2=-1$, set
\[
A=\mathcal C^+,
\]
and write every element uniquely as
\[
x=a+bI,
\qquad a,b\in A.
\]
The standard Clifford product uses the passive central law
\[
Ia=aI.
\]
The twin product keeps the same body, the same even core, and the same volume
line, but replaces that passive law by
\[
Ia=\widetilde a I.
\]
Thus the operation is not a random vector-space transport. It is an internal
deformation of the role of the volume element:
\[
\boxed{
\text{central volume law } Ia=aI
\quad\leadsto\quad
\text{conjugating volume law } Ia=\widetilde a I.
}
\]

This also clarifies the relation with generalized Clifford regularity. A
regularity theory should not depend only on a presented pair $(V,Q)$, nor only
on a naked algebra $\mathcal C$. It should specify the structured body and the
product/operations relative to which regularity is being defined. A natural
package is therefore something like
\[
(\mathcal C,\alpha,\widetilde{(\cdot)},\mathfrak v,\mu,D,G),
\]
where $\mu$ is the chosen multiplication, $D$ is the corresponding Dirac or
Cauchy--Riemann-type operator, and $G$ is the symmetry group being preserved.
The same Clifford body may support
\[
\mu_{\Cl}
\]
and
\[
\mu_{\OO},
\]
and these lead to different geometries, different spinor interpretations, and
potentially different regularity theories.

So the slogan must be corrected as follows. We should not say: ``forget the
underlying vector space and keep only the algebra.'' That loses too much. We
should say:
\begin{quote}
Forget the particular presentation $\Cl(V,Q)$, but keep the intrinsic graded
Clifford body: grading, reversion, even core, and oriented volume line.
\end{quote}
With that structure retained, the octonionic twin of $\Cl(3,0)$ is not an
external imposition. It is the result of activating the intrinsic central
volume line.

\section{Physical reading: switching the role of volume}

In ordinary $\Cl(3,0)$, the pseudoscalar
\[
I=e_1e_2e_3
\]
represents oriented volume. It distinguishes right-handed from left-handed
orientation, but it is passive. For a rotor $R\in H$,
\[
IR=RI.
\]
The volume element is therefore a background orientation marker, almost like a
global imaginary phase.

In the octonionic product, the same volume slot becomes active:
\[
IR=\bar R I=R^{-1}I
\]
for unit rotors $R$. Passing through the volume direction reverses the rotor.
The volume element is no longer merely the volume of the room; it becomes a
handedness-changing interface.

Thus the physical switch is:
\[
\Cl(3,0):\quad \text{space has an orientation,}
\]
whereas
\[
\OO:\quad \text{orientation acts on processes.}
\]
The nonassociativity records the time or position at which this chirality gate
is crossed. If $R_1$ and $R_2$ are noncommuting rotors, then
\[
I(R_1R_2)=\overline{R_1R_2}I=\bar R_2\bar R_1 I,
\]
while crossing the gate after only the first rotor produces the opposite order:
\[
(IR_1)R_2\sim \bar R_1\bar R_2 I.
\]
Since generally
\[
\bar R_1\bar R_2\neq \bar R_2\bar R_1,
\]
the two histories differ. Therefore
\[
\boxed{
\text{octonionic nonassociativity is memory of where chirality reversal occurred in a rotational process.}
}
\]

\section{Spinors in the two products}

The spinor interpretation makes the contrast especially transparent.

In ordinary $\Cl(3,0)$, spinors are the unit elements of the even subalgebra
\[
H=\Cl^+(3,0)\cong\mathbb H.
\]
A unit spinor or rotor $R\in H$ acts on a vector $v\in\mathbb R^3$ by
\[
v\mapsto RvR^{-1}.
\]
Thus
\[
\mathrm{Spin}(3)\cong SU(2)\cong \{R\in H:R\bar R=1\}.
\]
A Clifford spinor is therefore a rotational amplitude in a space with passive
orientation. The pseudoscalar $I$ commutes with it:
\[
IR=RI.
\]
The spinor and the oriented volume coexist without interfering. Composition is
ordinary group composition.

In the octonionic product, an element has the same linear form
\[
\Psi=a+bI,
\qquad a,b\in H.
\]
But now the second component is not merely another passive quaternionic copy.
It is the component that has crossed the active volume gate. Since
\[
IR=\bar R I=R^{-1}I,
\]
the $I$-component carries rotor information in conjugated, or chirality-flipped,
form.

Thus an octonionic spinor-like element may be read as
\[
\Psi=a+bI
\]
with
\[
a=\text{direct rotor/spinor channel},
\]
and
\[
bI=\text{chirality-gated rotor/spinor channel}.
\]
This is not an ordinary group spinor. Because the product is nonassociative,
composition depends on where the chirality gate enters the process.

For instance, for noncommuting rotors $R_1,R_2\in H$,
\[
(IR_1)R_2\neq I(R_1R_2)
\]
in general. The octonionic spinor therefore remembers whether the frame has
crossed the handedness gate before or after a given rotation.

In compressed form:
\[
\boxed{
\Cl(3,0):\quad \text{spinor = rotor in passive oriented 3-space.}
}
\]
\[
\boxed{
\OO:\quad \text{spinor-like element = rotor plus chirality-gated rotor.}
}
\]
Or, more physically:
\[
\boxed{
\text{Clifford spinors rotate frames; octonionic spinors remember whether a frame crossed the handedness gate.}
}
\]



\section{Folded and unfolded triality}

The discussion above also clarifies the relation with triality. The usual
route to triality starts from the large associative Clifford envelope
\[
\Cl(8,0),
\]
inside which
\[
\mathrm{Spin}(8)\subset \Cl^+(8,0)
\]
acts on three eight-dimensional representations:
\[
V,
\qquad
S_+,
\qquad
S_-.
\]
These are the vector representation and the two half-spinor representations.
In this presentation the three roles occupy different structural rooms. The
triality symmetry permutes them, but the Clifford architecture itself is large
and unfolded:
\[
\dim \Cl(8,0)=2^8=256.
\]

The octonionic presentation is different. One identifies
\[
V\cong \mathbb O,
\qquad
S_+\cong \mathbb O,
\qquad
S_-\cong \mathbb O.
\]
The three slots remain distinct as roles, but they are all modeled by the same
kind of eight-dimensional nonassociative body. The triality form may be written
schematically as
\[
t(x,\psi,\phi)=\operatorname{Re}(x(\psi\phi)),
\qquad
x,\psi,\phi\in\mathbb O,
\]
with the usual care about parenthesization. In this form triality is not spread
across a $256$-dimensional associative algebra. It is compressed into the
multiplicative geometry of $\mathbb O$.

This gives a useful distinction:
\[
\Cl(8,0)=\text{unfolded triality},
\]
while
\[
\mathbb O=\text{folded triality}.
\]
The word ``folded'' should not be read as a claim that the group
$\mathrm{Spin}(8)$ literally lives inside the ordinary $\Cl(3,0)$ product. It
does not. Rather, the active-volume twin of $\Cl(3,0)$ produces the octonionic
arena in which the three triality roles can be represented compactly.

Thus there are two routes:
\[
\Cl(8,0)
\longrightarrow
(V,S_+,S_-)
\longrightarrow
\text{triality},
\]
and
\[
\Cl(3,0)
\xrightarrow{\text{activate volume}}
\mathbb O
\longrightarrow
(\mathbb O_V,\mathbb O_{S_+},\mathbb O_{S_-}).
\]
The first route is the large associative envelope. The second route is the
minimal nonassociative triality cell.

This is the conceptual compression. In the unfolded Clifford picture, the
bookkeeping of triality is external: vector and spinor spaces are separate
modules. In the folded octonionic picture, that bookkeeping is internalized in
the nonassociative product. Parenthesization matters because the multiplication
must remember which triality face or channel is being used. In this sense:
\begin{quote}
Nonassociativity is compressed triality bookkeeping.
\end{quote}

The $\Cl(3,0)$ twin construction is therefore not merely an octonion trick. It
is a way of seeing how ordinary three-dimensional rotor geometry, once its
volume element is made active rather than passive, folds into the octonionic
cell that can carry vector, left-spinor, and right-spinor faces.

\section{The nonassociative twin principle}

We can now state the principle.

\begin{definition}[Clifford twin, informal]
A nonassociative twin of a Clifford algebra $\Cl(p,q)$ is a real
$2^n$-dimensional algebra, $n=p+q$, with basis indexed by $(\Ztwo)^n$,
multiplication
\[
u_a u_b = F(a,b)u_{a+b},
\]
and a nontrivial associator
\[
\assoc_F(a,b,c)\neq 1
\]
for some triples, while preserving as much as possible of the Clifford
sign geometry, grading, and quadratic signature.
\end{definition}

In this language:
\[
\Cl(3,0)
\quad\text{and}\quad
\OO
\]
are not the same algebra. But they are sibling algebras on the same
$8$-dimensional hypercube skeleton.

\begin{quote}
Clifford algebras are associative twisted hypercubes.
Octonions are nonassociative twisted hypercubes.
\end{quote}

There are two complementary ways to read this. At the combinatorial level:

\begin{quote}
The Clifford product twists signs while preserving associativity.
The octonion product twists signs so that associativity itself becomes curved.
\end{quote}

At the invariant $\Cl(3,0)$ level:

\begin{quote}
$\Cl(3,0)$ doubles $\mathbb H$ by a central volume imaginary.
The octonions double $\mathbb H$ by a conjugating volume-like unit.
\end{quote}



\section{Relation to prior work and the intended contribution}

The construction above must be positioned carefully. It should not be presented
as a new discovery of the octonions, nor as a claim that octonion multiplication
is simply Clifford multiplication in disguise. The basic ingredients are already
part of several well-developed literatures.

First, the Cayley--Dickson construction is the classical source of the
 octonion product. In modern notation,
\[
\mathbb O=\mathbb H\oplus \mathbb H\ell,
\]
with multiplication controlled by quaternionic conjugation. The formula used in
this note is exactly this construction, written in the internal language of
$\Cl(3,0)$ after identifying
\[
\Cl^+(3,0)\cong \mathbb H.
\]
Thus the present note does not claim to construct $\mathbb O$ for the first
time. It reinterprets the Cayley--Dickson step using the volume element of
three-dimensional Clifford geometry.

Second, the relationship between octonions, Clifford algebras, spinors, Bott
periodicity, triality, and exceptional structures is a classical theme. Baez's
survey \cite{BaezOctonions} is a standard reference for this landscape, and
Lounesto \cite{Lounesto} is a standard Clifford-algebra reference for the
spinorial side.

The same caution applies to triality. The representation-theoretic fact that
$\mathrm{Spin}(8)$ has vector and two half-spinor representations of dimension
$8$, cyclically related by triality, is classical. The octonionic description of
these representations and their triality form is also classical; it is treated
in Baez's account of spinors and triality and in the octonionic Clifford
representations of Manogue and Schray \cite{BaezOctonions,ManogueSchray}. The
present note uses this material only to articulate a contrast: $\Cl(8,0)$ is a
large associative envelope for unfolded triality, whereas $\mathbb O$ is the
compact nonassociative arena of folded triality.

Third, the twisted group algebra viewpoint is already known. Albuquerque and
Majid showed that the octonions arise as a quasiassociative twisting of the
group algebra of $(\Ztwo)^3$ \cite{AlbuquerqueMajidOctonions}. They later
showed that Clifford algebras may be obtained by twisting group algebras of
$(\Ztwo)^n$ by cocycles \cite{AlbuquerqueMajidClifford}. In that language,
the sharp algebraic contrast is that the Clifford twist is associative, while
the octonionic twist has a nontrivial associator. Basak gives another direct
treatment of the octonions as a twisted group algebra \cite{BasakOctonions}.
Morier--Genoud and Ovsienko develop related higher octonion-like families and
higher-dimensional analogues \cite{MorierGenoudOvsienko}. Elduque's discussion
of Clifford algebras as twisted group algebras and the Arf invariant is also
relevant to this background \cite{ElduqueArf}.

Fourth, there is recent work explicitly describing octonions as Clifford-like
nonassociative algebras. Depies, Smith, and Ashburn introduce Kingdon algebras:
alternative Clifford-like algebras equipped with an involutory anti-automorphism
and anticommuting generators, in which the octonions and split octonions arise
naturally \cite{DepiesSmithAshburn}. This is especially close in spirit to the
present note and must be treated as a central comparator in any future expanded
version.

The following table states the relation between the present proposal and these
nearby frameworks.

\begin{center}
\renewcommand{\arraystretch}{1.25}
\footnotesize
\begin{tabular}{>{\raggedright\arraybackslash}p{0.18\textwidth}
                >{\raggedright\arraybackslash}p{0.19\textwidth}
                >{\raggedright\arraybackslash}p{0.22\textwidth}
                >{\raggedright\arraybackslash}p{0.16\textwidth}
                >{\raggedright\arraybackslash}p{0.17\textwidth}}
\toprule
\textbf{Framework} & \textbf{Basic space} & \textbf{Control datum} & \textbf{Main identity / defect} & \textbf{Status here} \\
\midrule
Clifford algebra $\Cl(p,q)$ & Blade basis indexed by $(\Ztwo)^n$ & Quadratic form plus associative Clifford sign cocycle & Associativity; anticommutation of generators & Prior work / starting point \\
\addlinespace
Cayley--Dickson double & $A\oplus A\ell$ & Involution on $A$ and scalar $\ell^2$ & Associativity fails after the quaternion stage; alternativity survives for octonions & Prior work; formula reused \\
\addlinespace
Octonions $\mathbb O$ & $\mathbb H\oplus\mathbb H\ell$ or $(\Ztwo)^3$ basis & Quaternionic conjugation / Fano-plane sign rule & Nonassociative but alternative & Prior work; dimension-$8$ target \\
\addlinespace
Albuquerque--Majid twisted/quasi group algebras & Group algebra of $(\Ztwo)^n$ & Cocycle or quasi-cocycle / associator & Clifford: trivial associator; octonions: nontrivial associator & Prior work; algebraic comparator \\
\addlinespace
Morier--Genoud--Ovsienko higher octonions & Twisted group algebras over $(\Ztwo)^n$ & Cubic twisting functions & Higher octonion-like nonassociative laws & Prior work; higher-dimensional comparator \\
\addlinespace
Kingdon algebras & Clifford-like generators over formed spaces & Involutory anti-automorphism and alternative universal law & Alternative Clifford-like nonassociative structure & Prior work; closest conceptual comparator \\
\addlinespace
$\Cl(8,0)$ triality & $V,S_+,S_-$ inside a large associative envelope & $\mathrm{Spin}(8)$ action and outer triality symmetry & Three distinct $8$D representation roles & Prior work; unfolded triality comparator \\
\addlinespace
Octonionic triality & Three copies of $\mathbb O$ & Octonionic product and real trilinear form & One $8$D body occupies three triality roles & Prior work; folded triality comparator \\
\addlinespace
Intrinsic Clifford body viewpoint & $(\mathcal C,\alpha,\widetilde{(\cdot)},\mathfrak v)$ & Grading, reversion, and central volume line retained after forgetting a specific presentation & Product law becomes a variable part of the structure & Proposed contribution: invariant formulation \\
\addlinespace
Volume-twin interpretation in this note & $\Cl^+(3,0)\oplus\Cl^+(3,0)I$ & Reversion on $\Cl^+(3,0)$ plus active pseudoscalar $I$ & Associator records where chirality reversal occurs & Proposed contribution: geometric/physical reframing \\
\bottomrule
\end{tabular}
\end{center}

The invariant-body viewpoint adds one further clarification. The present note
is not merely changing a multiplication table on an anonymous $8$-dimensional
space. It keeps the graded Clifford body
\[
(\mathcal C,\alpha,\widetilde{(\cdot)},\mathfrak v)
\]
and changes the role of the central volume line inside that body. This is the
specific conceptual contribution: the Cayley--Dickson step is read as a
transition from passive central volume to active conjugating volume.

Thus the honest novelty claim is modest. We do not claim:
\begin{itemize}
    \item that the octonions are new;
    \item that the Cayley--Dickson construction is new;
    \item that octonions as twisted group algebras are new;
    \item that octonions as Clifford-like algebras are new;
    \item or that $\mathbb O$ is isomorphic to $\Cl(3,0)$ as an algebra.
\end{itemize}

The intended contribution is instead the following geometric reframing, together with the folded/unfolded triality interpretation:
\begin{quote}
Inside $\Cl(3,0)$, the vector-space decomposition
\[
\Cl(3,0)=\Cl^+(3,0)\oplus \Cl^+(3,0)I
\]
is already the Cayley--Dickson doubling space
\[
\mathbb H\oplus \mathbb H I.
\]
The Clifford product treats $I$ as a passive central volume imaginary. The
octonion product is obtained by making the same volume slot act as a
conjugating unit.
\end{quote}

In one line:
\[
\text{Clifford: } Ia=aI,
\qquad
\text{octonionic twin: } Ia=\bar a I.
\]

This is not intended as a replacement for the Cayley--Dickson construction. It
is an interpretation of the Cayley--Dickson construction from within the
geometry of $\Cl(3,0)$. The focus is on the change in the role of the oriented
volume element:
\[
\text{passive volume marker}
\quad\longrightarrow\quad
\text{active chirality gate}.
\]

This interpretation also gives a compact physical reading of the associator.
In the Clifford product, oriented volume is background structure. In the
octonionic product, crossing the volume direction conjugates a rotor, and hence
reverses the order of a composite rotation. Nonassociativity then records where
this chirality reversal occurs in a rotational process.

The possible research value is therefore not in rediscovering $\mathbb O$, but
in using this formulation as a template for a controlled comparison between:
\begin{enumerate}
    \item associative Clifford products;
    \item quasiassociative/twisted Boolean-cube products;
    \item Cayley--Dickson doubles of even Clifford cores;
    \item Kingdon-style alternative Clifford-like algebras;
    \item and computational searches for odd-dimensional volume twins.
\end{enumerate}

This is the exact level at which the present proposal should be defended: not
as a new algebraic object in dimension $8$, but as a useful invariant and
physical explanation of how the octonionic product can be placed on the same
underlying linear space as $\Cl(3,0)$ while changing the role of the oriented
volume element.

\section{A small research program}

The preceding discussion suggests a concrete program.

\begin{problem}[Classification of Clifford twins]
For each signature $(p,q)$, classify twists
\[
F:(\Ztwo)^n\times(\Ztwo)^n\to \{\pm 1\}
\]
such that:
\begin{enumerate}
    \item $A_F$ has dimension $2^n$;
    \item $A_F$ has the same grading skeleton as $\Cl(p,q)$;
    \item the square signs $u_a^2$ encode a prescribed quadratic signature;
    \item the associator $\assoc_F$ is nontrivial but controlled;
    \item optional identities hold: alternativity, flexibility, or power-associativity.
\end{enumerate}
\end{problem}

The dimension-$8$ case is the prototype. Higher dimensions should not be
expected to preserve all octonionic miracles. In particular, normed division
algebra behavior is exceptional. But the twisted-hypercube viewpoint gives
a uniform search space for Clifford-like nonassociative cousins.

\section{Computational viewpoint}

The computational representation is direct.

Represent elements of $(\Ztwo)^n$ as bit vectors or integers
\[
0,\dots,2^n-1.
\]
A twist is a table
\[
F[a,b]\in\{\pm 1\}.
\]
Multiplication is
\[
(a,b)\mapsto (F[a,b],a\oplus b),
\]
where $\oplus$ denotes bitwise XOR.

The associator table is
\[
\assoc[a,b,c]
=
\frac{F[a,b]F[a\oplus b,c]}
     {F[b,c]F[a,b\oplus c]}.
\]
This makes the theory ideal for a small Python package:
\begin{itemize}
    \item generate Clifford twists;
    \item generate Fano/octonion twists;
    \item compute associators;
    \item test algebraic identities;
    \item search for nonassociative twins under constraints.
\end{itemize}


\section{Conclusion}

The slogan ``Clifford algebras have nonassociative twins'' is mathematically
reasonable only when stated carefully. In the Boolean-cube language, Clifford
algebras and octonionic-type algebras share a graded hypercube skeleton, but
differ in whether the associator is trivial or nontrivial. This is the broad
combinatorial view.

In dimension $8$ there is a stronger geometric view. The Clifford algebra
$\Cl(3,0)$ contains an intrinsic quaternionic even subalgebra
\[
\Cl^+(3,0)\cong\mathbb H,
\]
and an intrinsic split
\[
\Cl(3,0)=\Cl^+(3,0)\oplus\Cl^+(3,0)I.
\]
The Clifford product treats $I$ as a passive central volume imaginary. The
octonion product is obtained by keeping the same rotor core and the same
doubled vector space, but replacing the central rule
\[
Ia=aI
\]
by the conjugating rule
\[
\ell a=\bar a\ell.
\]
Thus the octonionic twin of $\Cl(3,0)$ is not merely a basis-level sign trick.
It is the result of turning the passive volume direction into an active
handedness-reversing gate for the rotor algebra.

The Cayley--Dickson unit is also not arbitrary if rotational invariance is
required. The odd part of $\Cl(3,0)$ contains vectors, which select directions,
and the pseudoscalar line $\mathbb RI$, which selects only orientation. Hence
$I$ is the unique isotropic candidate, up to sign. In this sense the
octonionic extension is chiral rather than directional.

The corresponding spinor picture is equally sharp: Clifford spinors are rotors
in passively oriented three-space, while octonionic spinor-like elements are
rotors together with chirality-gated rotors. Their nonassociativity records
where the handedness gate occurs in a rotational process.

The construction should therefore be formulated not from the naked algebra
alone, and not merely from the presentation $\Cl(V,Q)$, but from an oriented
graded Clifford body
\[
(\mathcal C,\alpha,\widetilde{(\cdot)},\mathfrak v).
\]
For $\Cl(3,0)$ the center recovers the volume line $\mathbb RI$, while the
orientation chooses $I$ or $-I$. The grade involution recovers the even core,
and reversion gives the conjugation needed for the Cayley--Dickson step. In
this precise sense, the octonionic twin is an intrinsic volume activation of
$\Cl(3,0)$, not an arbitrary remultiplication of $\mathbb R^8$.

Finally, the same construction gives a compact way to think about triality.
The usual $\Cl(8,0)$ route unfolds triality across the vector and two half-spinor
representations inside a large associative envelope. The octonionic route folds
those three roles into one nonassociative eight-dimensional body. The
$\Cl(3,0)$ twin construction does not replace $\mathrm{Spin}(8)$ and does not
place full triality inside the ordinary Clifford product. What it does is
produce, from three-dimensional rotor geometry by active volume, the octonionic
cell in which folded triality naturally lives.

This suggests a focused research direction: use the invariant Clifford-body
viewpoint to compare associative Clifford products, volume-activated twins,
twisted Boolean-cube models, and octonionic triality cells, while remaining
explicit about what is prior work and what is only a new geometric synthesis.

\begin{thebibliography}{99}

\bibitem{CayleyDicksonBackground}
R. D. Schafer,
\emph{An Introduction to Nonassociative Algebras},
Academic Press, 1966.

\bibitem{SpringerVeldkamp}
T. A. Springer and F. D. Veldkamp,
\emph{Octonions, Jordan Algebras and Exceptional Groups},
Springer Monographs in Mathematics, Springer, 2000.

\bibitem{BaezOctonions}
J. C. Baez,
\emph{The Octonions},
Bulletin of the American Mathematical Society 39, no. 2 (2002), 145--205.
DOI: \url{https://doi.org/10.1090/S0273-0979-01-00934-X}.

\bibitem{Lounesto}
P. Lounesto,
\emph{Clifford Algebras and Spinors},
2nd edition, London Mathematical Society Lecture Note Series 286,
Cambridge University Press, 2001.

\bibitem{AlbuquerqueMajidOctonions}
H. Albuquerque and S. Majid,
\emph{Quasialgebra structure of the octonions},
Journal of Algebra 220, no. 1 (1999), 188--224.
Available as arXiv:math/9802116.

\bibitem{AlbuquerqueMajidClifford}
H. Albuquerque and S. Majid,
\emph{Clifford algebras obtained by twisting of group algebras},
Journal of Pure and Applied Algebra 171, no. 2--3 (2002), 133--148.
Available as arXiv:math/0011040.

\bibitem{BasakOctonions}
T. Basak,
\emph{The octonions as a twisted group algebra},
arXiv:1702.05705, 2017.

\bibitem{MorierGenoudOvsienko}
S. Morier-Genoud and V. Ovsienko,
\emph{A series of algebras generalizing the octonions and Hurwitz--Radon identity},
Communications in Mathematical Physics 306, no. 1 (2011), 83--118.
Available as arXiv:1003.0429.

\bibitem{ElduqueArf}
A. Elduque,
\emph{Clifford algebras as twisted group algebras and the Arf invariant},
Advances in Applied Clifford Algebras 28 (2018), article 42.

\bibitem{DepiesSmithAshburn}
C. M. Depies, J. D. H. Smith, and M. D. Ashburn,
\emph{Octonions as Clifford-like algebras},
arXiv:2310.09972, 2023.

\end{thebibliography}

\end{document}
